% !TEX program = xelatex
% NOTE: 为了更好的 Unicode 支持，请使用 XeLaTeX（或 LuaLaTeX）编译。
\documentclass[11pt,aspectratio=169]{beamer}
\usetheme{Madrid}
\usecolortheme{default}
\usefonttheme{professionalfonts}
\setbeamertemplate{navigation symbols}{}
\setbeamertemplate{caption}[numbered]
\usepackage{fontspec}
\usepackage{xeCJK}
\setmainfont{CMU Serif}
\IfFontExistsTF{Noto Serif CJK SC}{\setCJKmainfont{Noto Serif CJK SC}}{%
  \IfFontExistsTF{Source Han Serif SC}{\setCJKmainfont{Source Han Serif SC}}{%
    \setCJKmainfont{FandolSong-Regular}% TeX Live fallback
  }%
}%
\usepackage{amsmath}
\usepackage{graphicx}
\usepackage{tikz-cd}
\usepackage{hyperref}

\begin{document}

\title[ACT-FedViM]{ACT-FedViM：面向海洋浮游生物监测的自适应联邦子空间学习与 OOD 检测框架}
\author{孙伟洲}
\institute{北京师范大学珠海校区}
\date{\today}
\frame{\titlepage}

\begin{frame}{目录}
  \tableofcontents
\end{frame}

\section{研究思路与核心框架 (Idea \& Framework)}



\begin{frame}{ACT-FedViM}
  \centering
  \vfill
  {\LARGE 研究思路与核心框架\\[0.5em]}
  {\large Idea \& Framework}
  \vfill
\end{frame}

\begin{frame}{研究思路}
  \begin{block}{ACT-FedViM 解决的三大核心挑战}
    \begin{itemize}
      \item 针对海洋监测节点地理分布广泛且通信带宽受限的挑战：采用联邦学习 (FL) 框架，避免海量图像数据直接传输。
      \item 针对边缘设备计算资源有限的挑战：针对联邦场景对 ViM 进行统计量驱动的适配，采用基于统计残差的 ViM-style 后处理算法，即插即用，摒弃高开销的生成式密度估计框架。
      \item 针对深度学习分类模型的高维特征冗余与 PCA 主子空间估计不稳问题：引入 ACT (Adjusted Correlation Thresholding) 机制，自动截断噪声谱尾，为主子空间规模选择提供更有统计依据的自适应规则。
    \end{itemize}
  \end{block}
\end{frame}

\begin{frame}{算法整体流程 (Algorithm Pipeline)}
  \centering
  \includegraphics[width=\linewidth]{pipline.png}
\end{frame}

\section{实验配置与结果 (Experimental Results)}

\begin{frame}{ACT-FedViM}
  \centering
  \vfill
  {\LARGE 实验配置与结果\\[0.5em]}
  {\large Experimental Results}
  \vfill
\end{frame}



\begin{frame}{实验配置与结果}
  \small
  \begin{enumerate}
    \item \textbf{DYB 数据集与 Non-IID 设定 (Federated Non-IID Setup)}
      \begin{itemize}
        \item 基础数据集：长尾细粒度海洋浮游生物数据集 (DYB-PlanktonNet)。
        \item 分布内 (ID) 规模：包含 54 个已知 ID 类别；约 26034 张图像用于联邦训练，2603 张图像作为验证集，2939 张作为测试集，用于分类准确率评估。
        \item 联邦学习数据模拟：采用狄利克雷分布 (Dirichlet Distribution, $\alpha=0.1$) 将训练数据划分至 5 Clients。
      \end{itemize}

    \item \textbf{OOD 评测基准 (OOD Datasets Benchmark)}
      \begin{itemize}
        \item 评测阶段引入 18547 张分布外图像，划分为两部分，分别计算其对应的 AUROC：
        \item 近分布外 (Near-OOD)：包含 25 类，共 1516 张；由在形态或生态属性上与 ID 类别相关，但不属于训练目标类别的浮游生物构成。
        \item 远分布外 (Far-OOD)：包含 12 类，共 17031 张；由与 ID 类别在语义上显著偏离的样本构成，选取域内非生物或成像干扰项（如气泡、颗粒、碎片等）。
      \end{itemize}
    
  \end{enumerate}
\end{frame}

\begin{frame}{ID vs. OOD 示例}
  \centering
  \includegraphics[height=0.78\textheight,keepaspectratio]{IDandOOD.png}
\end{frame}

\begin{frame}{实验配置结果（续）}
  \small
  \begin{enumerate}
    \setcounter{enumi}{2}
    \item \textbf{参评的九大主流视觉网络 (9 Backbone Architectures)}
      \begin{itemize}
        \item 经典 CNN：ResNet50, ResNet101, DenseNet169
        \item 现代 CNN：EfficientNetV2-S, MobileNetV3-Large, ConvNeXt-Base
        \item Transformer：ViT-B/16, ViT-B/32, DeiT-Base
      \end{itemize}

    \item \textbf{模型结果总结}
      \begin{itemize}
        \item 采用保留前 95\% 的方差贡献率筛选主成分个数的方法，所有模型均展现出优异的 OOD 检测能力：ID 分类准确率在 95.37\%--96.97\% 之间，Near-OOD AUROC 在 87.23\%--96.64\% 之间，Far-OOD AUROC 在 88.18\%--97.58\% 之间。
        \item CNN 类模型与 Transformer 类模型在原版 ViM 下性能差距不大，说明 ViM 方法对不同架构具有良好的通用性。
        \item 应用 ACT（Adjusted Correlation Thresholding）算法自动筛选子空间维度后，模型表现呈现明显的架构依赖性：CNN 模型平均压缩率为 78.6\%，Near-OOD AUROC 平均提升 +3.61\%；Transformer 模型平均压缩率为 74.9\%，但性能变化仅为 ±0.19\%，基本保持稳定。
        \item 差异原因：CNN 特征空间呈现陡峭的特征值衰减（主成分集中），而 Transformer 特征值分布平缓（信息分散），导致 ACT 对前者降噪效果显著，对后者作用有限。
      \end{itemize}
  \end{enumerate}
\end{frame}

\begin{frame}{不同Backbone的性能对比}
  \centering
  \includegraphics[height=0.82\textheight,keepaspectratio]{9model_performance_histogram.png}
\end{frame}

\begin{frame}[fragile]{不同Backbone的性能对比}
  \scriptsize
  \setlength{\tabcolsep}{4pt}
  \renewcommand{\arraystretch}{1.15}
  \begin{tabular}{lcccc}
    \hline
    模型名称 &
    (a) ID acc (\%) &
    (b) Near-OOD AUROC (\%)  &
    (c) Far-OOD AUROC (\%)  &
    (d) 空间压缩率 (\%) \\
    \hline
    ResNet101 & 96.97\% & 93.7 / 96.9 (+3.29\%) & 94.5 / 96.9 (+2.35\%) & 75.3\% \\
    ResNet50 & 96.60\% & 93.8 / 96.7 (+2.83\%) & 90.8 / 95.6 (+4.74\%) & 75.3\% \\
    DenseNet169 & 95.92\% & 91.1 / 95.2 (+4.08\%) & 92.3 / 96.7 (+4.36\%) & 76.0\% \\
    EfficientNetV2-S & 96.67\% & 87.2 / 96.0 (+8.80\%) & 88.2 / 96.5 (+8.31\%) & 92.6\% \\
    ConvNeXt-Base & 96.26\% & 96.7 / 97.0 (+0.28\%) & 97.6 / 96.5 (-1.06\%) & 61.2\% \\
    MobileNetV3-L & 95.54\% & 93.8 / 96.1 (+2.35\%) & 95.7 / 97.4 (+1.63\%) & 78.7\% \\
    ViT-B/16 & 95.68\% & 96.1 / 96.0 (-0.07\%) & 94.8 / 94.8 (+0.01\%) & 72.5\% \\
    ViT-B/32 & 95.37\% & 95.6 / 95.7 (+0.05\%) & 94.4 / 94.2 (-0.15\%) & 73.3\% \\
    DeiT-Base & 96.09\% & 96.5 / 96.0 (-0.54\%) & 96.7 / 96.3 (-0.39\%) & 78.2\% \\
    \hline
  \end{tabular}

  \vspace{0.6em}
  {\scriptsize 各个模型的参数量从大到小分别是：ConvNeXt-Base 87.62M、ViT-B/32 87.50M、ViT-B/16 85.84M、DeiT-Base 85.84M、ResNet101 42.61M、ResNet50 23.62M、EfficientNetV2-S 20.25M、DenseNet169 12.57M、MobileNetV3-L 3.02M}
\end{frame}

\section{关键技术原理 (Technical Breakdown)}

\begin{frame}{ACT-FedViM}
  \centering
  \vfill
  {\LARGE 关键技术原理\\[0.5em]}
  {\large Technical Breakdown}
  \vfill
\end{frame}


\begin{frame}{1. 联邦学习的全局协方差重构}
  \small

  \hspace{2em}在高度异构（Non-IID）设定下，单个客户端的局部特征子空间会发生偏斜，无法代表真实的数据流形。为了能够还原整体数据的协方差矩阵，本框架采用\textbf{充分统计量（Sufficient Statistics）}传递机制。

  \hspace{2em}对于 $K$ 个客户端，设第 $k$ 个客户端拥有 $n_k$ 个样本，特征提取器 $f_\theta$ 输出的特征为 $z \in \mathbb{R}^D$。在模型训练收敛后，各客户端仅计算并上传本地特征分布的低阶矩：
  \begin{itemize}
    \item 零阶矩（样本量）：$n_k$
    \item 一阶矩（特征和）：$S_k^{(1)} = \sum z$
    \item 二阶矩（未中心化协方差和）：$S_k^{(2)} = \sum z z^\top$
  \end{itemize}

  \hspace{2em}服务端通过聚合局部统计量，利用了协方差的分解性质，在数学上无偏地重构出全局样本均值 $\mu_{\mathrm{global}}$ 和全局总体协方差矩阵 $\Sigma_{\mathrm{global}}$：
  \[
    \mu_{\mathrm{global}} = \frac{\sum_{k=1}^K S_k^{(1)}}{\sum_{k=1}^K n_k}
  \]
  \[
    \Sigma_{\mathrm{global}} =
    \frac{1}{N_{\mathrm{total}}-1}
    \left(
      \sum_{k=1}^K S_k^{(2)} - N_{\mathrm{total}} \cdot \mu_{\mathrm{global}}\mu_{\mathrm{global}}^\top
    \right)
  \]
\end{frame}

\begin{frame}[t]{2. 随机矩阵理论与 ACT 自适应降噪 (Random Matrix Theory \& ACT)}
  \small

  \hspace{2em}（1）获取 $\Sigma_{\mathrm{global}}$ 后，原始 ViM 方法需要人为设定主子空间维度 $k$，论文中主要采用固定维度并讨论其在一定范围内的稳健性。在本文的联邦化 ViM 基线实现中，我们进一步采用固定方差保留率（如 95\%）的 PCA 作为一个动态 heuristic 来确定原始子空间规模。然而，现代视觉骨干网络（如 ResNet50, DenseNet169）的末端特征维度 $p$ 极大，通常在 1024 到 2048 维之间，相应一个 $p \times p$ 的协方差矩阵，其需要估计的独立自由度（参数量）是 $\frac{p(p+1)}{2}$。相比之前用于训练的总样本量 $n$（如数万张浮游生物图像）虽然庞大，但在数量级仍然远小于$\frac{p(p+1)}{2}$。\par\vspace{0.4em}
\par \hspace{2em}（2）这种“大维度 $p$、大样本 $n$”的数据结构（即维度与样本量之比 $p/n \to \gamma > 0$）直接打破了经典统计学中“样本量远大于维度（$n \gg p$）”的渐进假设。在这一高维机制下，样本协方差矩阵的经验特征值将不再是真实总体特征值的无偏估计。\par\vspace{0.4em}
\par \hspace{2em}（3）Marchenko-Pastur 定律指出，高维纯噪声数据亦会产生显著的非零特征值，导致经典 PCA 将大量噪声维度误认为信号；这些谱尾噪声还会进一步干扰主子空间方向的稳定估计。\par\vspace{0.4em}
  \end{frame}
  
  
  
  
\begin{frame}[t]{2. 随机矩阵理论与 ACT 自适应降噪 (Random Matrix Theory \& ACT)}
  \small
\par 因此，本文引入 ACT 的核心动机不是否定 ViM 的子空间思想，而是沿着“原始 ViM 的固定 $k$ 设定 $\rightarrow$ 本文联邦基线的方差贡献率 heuristic $\rightarrow$ ACT 的统计驱动替代”这一路径，逐步提升主子空间维度选择的合理性。ACT 的作用是为本文联邦 ViM 基线中经验性的子空间规模 heuristic 提供一个更有统计依据的自适应替代，并缓解高维 PCA 下主子空间方向难以稳定估计的问题。为解决此高维统计学难题，本框架引入ACT算法，对特征空间进行降维去噪：
  \begin{itemize}
    \item \textbf{标准化与阈值界定：}将 $\Sigma_{\mathrm{global}}$ 转化为相关矩阵 $R$ 以消除特征尺度差异；计算零假设（纯高斯噪声）下的 Marchenko--Pastur 分布上界阈值
    \[
      s = 1 + \sqrt{\frac{p}{n-1}}.
    \]
    \item \textbf{高维偏差修正：}样本特征值 $\lambda_j$ 因维度效应发生膨胀。引入离散 Stieltjes 变换 $m_{nj}(z)$，结合解析修正项重构 $\underline{m}_{nj}(z)$，反解出真实的总体特征值估计：
    \[
      \lambda_j^C = -\frac{1}{\underline{m}_{nj}(\lambda_j)}.
    \]
    \item \textbf{自适应截断：}最终的主子空间维度 $k$ 自动确定为：
    \[
      k_{\mathrm{optimal}} = \max \{\, j : \lambda_j^C > s \,\}.
    \]
  \end{itemize}

\par ACT 算法精准切断了特征值谱的右侧长尾，并为本文联邦 ViM 基线中的经验性主子空间维度选择提供了一个统计驱动的替代规则；通过先去除噪声主导的谱尾，再执行 PCA，也能让最终主子空间方向的估计更加稳定。
\end{frame}

\begin{frame}{3. ViM 与校准系数$\alpha$}
  \small


  \hspace{2em}原版 ViM（Virtual-logit Matching）的核心思想是结合特征空间的几何约束与输出空间的概率置信度。本文针对联邦场景对其进行了统计量驱动的改造：使用联邦重构的全局均值 $\mu_{\mathrm{global}}$ 与全局子空间进行中心化和打分。对于输入样本 $x$，其特征为 $z$，分类层 Logits 向量为 $v$：
  \begin{itemize}
    \item \textbf{几何残差 (Residual)：}度量样本偏离已知数据主流形的距离，
    \[
      S_{\mathrm{res}}(x)=\left\|(I-PP^\top)(z-\mu_{\mathrm{global}})\right\|_2.
    \]
    \item \textbf{概率能量 (Energy)：}利用 LogSumExp 算子估计边缘对数概率，
    \[
      S_{\mathrm{energy}}(x)=\log\sum_{c=1}^C e^{v_c}.
    \]
  \end{itemize}

  \hspace{2em}\textbf{经验校准系数 $\alpha$：}几何残差与概率能量属于不同的量纲。本文在最终评估阶段直接利用 ID 训练特征进行经验校准：
  \[
    \mu_{\mathrm{res}} =
    \frac{1}{N}\sum_{i=1}^{N}\left\|(I-PP^\top)(z_i-\mu_{\mathrm{global}})\right\|_2.
  \]
  \[
    \mu_{\mathrm{energy}} =
    \frac{1}{N}\sum_{i=1}^{N}\mathrm{LogSumExp}(v_i).
  \]
  \[
    \alpha=\frac{|\mu_{\mathrm{energy}}|}{\mu_{\mathrm{res}}+\epsilon}.
  \]
\end{frame}



\begin{frame}{3. ViM 与校准系数$\alpha$}
  \small
  \hspace{2em}最终，测试样本的联合 OOD 得分为：
  \[
    \mathrm{Score}_{\mathrm{OOD}}(x)=S_{\mathrm{energy}}(x)-\alpha\cdot S_{\mathrm{res}}(x).
  \]
  \hspace{2em}分布内（ID）样本具备高能量与低残差，得分较高；而 OOD 样本具有低能量与高残差，经 $\alpha$ 对齐量纲后实现与 ID 样本的分离。
\end{frame}



\end{document}
